\olchapter{fol}{mat}{Theories and Their Models}



\olfileid{fol}{mat}{int}
\olsection{Introduction}

\begin{explain}
The development of the axiomatic method is a significant achievement
in the history of science, and is of special importance in the history
of mathematics.  An axiomatic development of a field involves the
clarification of many questions: What is the field about?  What are the
most fundamental concepts?  How are they related?  Can all the
concepts of the field be defined in terms of these fundamental
concepts?  What laws do, and must, these concepts obey?

The axiomatic method and logic were made for each other.  Formal logic
provides the tools for formulating axiomatic theories, for proving
theorems from the axioms of the theory in a precisely specified way,
for studying the properties of all systems satisfying the axioms in a
systematic way.
\end{explain}

\begin{defn}
A set of !!{sentence}s~$\Gamma$ is \emph{closed} iff, whenever
$\Gamma \Entails !A$ then $!A \in \Gamma$.  The \emph{closure} of a set
of !!{sentence}s~$\Gamma$ is $\Setabs{!A}{\Gamma \Entails !A}$.

We say that~$\Gamma$ is \emph{axiomatized by} a set of
sentences~$\Delta$ if $\Gamma$ is the closure of~$\Delta$.
\end{defn}

\begin{explain}
We can think of an axiomatic theory as the set of sentences that is
axiomatized by its set of axioms~$\Delta$.  In other words, when we
have a first-order language which contains non-logical symbols for the
primitives of the axiomatically developed science we wish to study,
together with a set of !!{sentence}s that express the fundamental laws
of the science, we can think of the theory as represented by all the
!!{sentence}s in this language that are entailed by the axioms.  This
ranges from simple examples with only a single primitive and
simple axioms, such as the theory of partial orders, to complex
theories such as Newtonian mechanics.

The important logical facts that make this formal approach to the
axiomatic method so important are the following.  Suppose $\Gamma$ is
an axiom system for a theory, i.e., a set of sentences.
\begin{enumerate}
\item We can state precisely when an axiom system captures an intended
  class of !!{structure}s.  That is, if we are interested in a certain
  class of !!{structure}s, we will successfully capture that class by
  an axiom system~$\Gamma$ iff the !!{structure}s are exactly
  those~$\Struct M$ such that $\Sat{M}{\Gamma}$.
\item We may fail in this respect because there are $\Struct M$ such
  that $\Sat{M}{\Gamma}$, but $\Struct M$ is not one of the
  !!{structure}s we intend.  This may lead us to add axioms which are
  not true in~$\Struct M$.
\item If we are successful at least in the respect that $\Gamma$ is
  true in all the intended !!{structure}s, then a sentence~$!A$ is true in
  all intended !!{structure}s whenever $\Gamma \Entails !A$.  Thus we can
  use logical tools (such as !!{derivation} methods) to show that sentences are
  true in all intended !!{structure}s simply by showing that they are
  entailed by the axioms.
\item Sometimes we don't have intended !!{structure}s in mind, but instead
  start from the axioms themselves: we begin with some primitives that
  we want to satisfy certain laws which we codify in an axiom system.
  One thing that we would like to verify right away is that the axioms
  do not contradict each other: if they do, there can be no concepts
  that obey these laws, and we have tried to set up an incoherent
  theory.  We can verify that this doesn't happen by finding a model
  of~$\Gamma$.  And if there are models of our theory, we can use
  logical methods to investigate them, and we can also use logical
  methods to construct models.
\item The independence of the axioms is likewise an important
  question.  It may happen that one of the axioms is actually a
  consequence of the others, and so is redundant.  We can prove that
  an axiom $!A$ in $\Gamma$ is redundant by proving $\Gamma \setminus
  \{!A\} \Entails !A$.  We can also prove that an axiom is not
  redundant by showing that $(\Gamma \setminus \{!A\}) \cup \{\lnot
  !A\}$ is satisfiable.  For instance, this is how it was shown that the
  parallel postulate is independent of the other axioms of geometry.
\item Another important question is that of definability of concepts
  in a theory: The choice of the language determines what the models
  of a theory consist of.  But not every aspect of a theory must be
  represented separately in its models.  For instance, every ordering
  $\le$ determines a corresponding strict ordering~$<$---given one, we
  can define the other.  So it is not necessary that a model of a
  theory involving such an order must \emph{also} contain the
  corresponding strict ordering.  When is it the case, in general,
  that one relation can be defined in terms of others?  When is it
  impossible to define a relation in terms of others (and hence must
  add it to the primitives of the language)?
\end{enumerate}
\end{explain}






\olfileid{fol}{mat}{exs}

\olsection{Expressing Properties of \printtoken{P}{structure}}

\begin{explain}
It is often useful and important to express conditions on
functions and relations, or more generally, that the functions and
relations in a structure satisfy these conditions.  For instance, we
would like to have ways of distinguishing those !!{structure}s for a
language which ``capture'' what we want the !!{predicate}s to ``mean''
from those that do not.  Of course we're completely free to specify
which !!{structure}s we ``intend,'' e.g., we can specify that the
interpretation of the !!{predicate}~$\le$ must be an ordering, or that
we are only interested in interpretations of~$\Lang L$ in which the
domain consists of sets and $\Obj \in$ is interpreted by the ``is
!!a{element} of'' relation.  But can we do this with !!{sentence}s of
the language?  In other words, which conditions on
!!a{structure}~$\Struct M$ can we express by !!a{sentence} (or perhaps
a set of !!{sentence}s) in the language of~$\Struct M$?  There are
some conditions that we will not be able to express.  For instance,
there is no sentence of~$\Lang L_A$ which is only true in a
!!{structure}~$\Struct M$ if $\Domain M = \Nat$.  We cannot express
``the domain contains only natural numbers.''  But there are
``structural properties'' of !!{structure}s that we perhaps can
express.  Which properties of !!{structure}s can we express by
!!{sentence}s?  Or, to put it another way, which collections of
!!{structure}s can we describe as those making !!a{sentence} (or set
of !!{sentence}s) true?
\end{explain}

\begin{defn}[Model of a set]
Let $\Gamma$ be a set of !!{sentence}s in a language~$\Lang L$.  We
say that !!a{structure}~$\Struct M$ \emph{is a model of}~$\Gamma$ if
$\Sat{M}{!A}$ for all $!A \in \Gamma$.
\end{defn}

\begin{ex}
The sentence $\lforall[x][x \le x]$ is true in~$\Struct M$ iff
$\Assign{\le}{M}$ is a reflexive relation.  The sentence
$\lforall[x][\lforall[y][((x \le y \land y \le x) \lif x = y)]]$ is
true in~$\Struct M$ iff $\Assign{\le}{M}$ is anti-symmetric.  The
sentence $\lforall[x][\lforall[y][\lforall[z][((x \le y \land y \le z)
      \lif x \le z)]]]$ is true in~$\Struct M$ iff $\Assign{\le}{M}$
is transitive.  Thus, the models of
\begin{align*}
\{\quad &\lforall[x][x \le x], \\
   & \lforall[x][\lforall[y][((x \le y \land y \le
    x) \lif x = y)]], \\
   &\lforall[x][\lforall[y][\lforall[z][((x \le y
      \land y \le z) \lif x \le z)]]] \quad \}
\end{align*}
are exactly those structures in which~$\Assign{\le}{M}$ is reflexive,
anti-symmetric, and transitive, i.e., a partial order.  Hence, we can take
them as axioms for the \emph{first-order theory of partial orders}.
\end{ex}





\olfileid{fol}{mat}{the}

\olsection{Examples of First-Order Theories}

\begin{ex}
The theory of strict linear orders in the language~$\Lang L_<$ is
axiomatized by the set
\begin{align*}
\{\quad & \lforall[x][\lnot x < x], \\
& \lforall[x][\lforall[y][((x < y \lor y <
    x) \lor x = y)]], \\
& \lforall[x][\lforall[y][\lforall[z][((x < y
      \land y < z) \lif x < z)]]] \quad \}
\end{align*}
It completely captures the intended !!{structure}s: every strict
linear order is a model of this axiom system, and vice versa, if $R$
is a linear order on a set $X$, then the structure $\Struct M$ with
$\Domain M = X$ and $\Assign{<}{M} = R$ is a model of this theory.
\end{ex}

\begin{ex}
The theory of groups in the language $\Obj 1$ (!!{constant}), $\cdot$
(two-place !!{function}) is axiomatized by
\begin{align*}
& \lforall[x][\eq[(x \cdot \Obj 1)][x]]\\
& \lforall[x][\lforall[y][\lforall[z][\eq[(x \cdot (y \cdot z))][((x
          \cdot y) \cdot z)]]]]\\
& \lforall[x][\lexists[y][\eq[(x \cdot y)][\Obj 1]]]
\end{align*}
\end{ex}

\begin{ex}
The theory of Peano arithmetic is axiomatized by the following
sentences in the language of arithmetic~$\Lang L_A$.
\begin{align*}
& \lforall[x][\lforall[y][(\eq[x'][y'] \lif \eq[x][y])]]\\
& \lforall[x][\eq/[\Obj 0][x']]\\
& \lforall[x][\eq[(x + \Obj 0)][x]]\\
& \lforall[x][\lforall[y][\eq[(x + y')][(x + y)']]]\\
& \lforall[x][\eq[(x \times \Obj 0)][\Obj 0]]\\
& \lforall[x][\lforall[y][\eq[(x \times y')][((x \times y) + x)]]]\\
& \lforall[x][\lforall[y][(x < y \liff \lexists[z][\eq[(z' + x)][y])]]]\\
\intertext{plus all sentences of the form}
& (!A(\Obj 0) \land \lforall[x][(!A(x) \lif !A(x'))]) \lif \lforall[x][!A(x)]
\end{align*}
Since there are infinitely many sentences of the latter form, this
axiom system is infinite.  The latter form is called the
\emph{induction schema}. (Actually, the induction schema is a bit more
complicated than we let on here.)

The last axiom is an \emph{explicit definition} of~$<$.
\end{ex}

\begin{ex}
The theory of pure sets plays an important role in the foundations
(and in the philosophy) of mathematics.  A set is pure if all its
!!{element}s are also pure sets.  The empty set counts therefore as
pure, but a set that has something as !!a{element} that is not a set
would not be pure.  So the pure sets are those that are formed just
from the empty set and no ``urelements,'' i.e., objects that are not
themselves sets.

The following might be considered as an axiom system for a theory of
pure sets:
\begin{align*}
& \lexists[x][\lnot \lexists[y][y \in x]]\\
& \lforall[x][\lforall[y][(\lforall[z](z \in x \liff z \in y) \lif 
\eq[x][y])]]\\
& \lforall[x][\lforall[y][\lexists[z][\lforall[u][(u \in z \liff
(\eq[u][x] \lor \eq[u][y]))]]]]\\
& \lforall[x][\lexists[y][\lforall[z][(z \in y \liff \lexists[u][(z \in
        u \land u \in x)])]]]\\
\intertext{plus all sentences of the form} &
\lexists[x][\lforall[y][(y \in x \liff !A(y))]]
\end{align*}
The first axiom says that there is a set with no !!{element}s (i.e.,
$\emptyset$ exists); the second says that sets are extensional; the
third that for any sets $X$ and $Y$, the set $\{X, Y\}$ exists; the
fourth that for any set $X$, the set $\cup X$ exists, where $\cup X$ is the 
union of all the elements of $X$.

The !!{sentence}s mentioned last are collectively called the
\emph{naive comprehension scheme}.  It essentially says that for every
$!A(x)$, the set $\Setabs{x}{!A(x)}$ exists---so at first glance a
true, useful, and perhaps even necessary axiom.  It is called ``naive''
because, as it turns out, it makes this theory unsatisfiable: if you
take $!A(y)$ to be $\lnot y \in y$, you get the !!{sentence}
\[
\lexists[x][\lforall[y][(y \in x \liff \lnot y \in y)]]
\]
and this !!{sentence} is not satisfied in any !!{structure}.
\end{ex}

\begin{ex}
In the area of \emph{mereology}, the relation of \emph{parthood} is a
fundamental relation.  Just like theories of sets, there are theories
of parthood that axiomatize various conceptions (sometimes
conflicting) of this relation.

The language of mereology contains a single two-place predicate
symbol~$\Obj P$, and $\Atom{\Obj P}{x, y}$ ``means'' that $x$ is a
part of~$y$.  When we have this interpretation in mind, !!a{structure}
for this language is called a \emph{parthood structure}.  Of course,
not every structure for a single two-place predicate will really
deserve this name.  To have a chance of capturing ``parthood,''
$\Assign{\Obj P}{M}$ must satisfy some conditions, which we can lay
down as axioms for a theory of parthood.  For instance, parthood is a
partial order on objects: every object is a part (albeit an
\emph{improper} part) of itself; no two different objects can be parts
of each other; a part of a part of an object is itself part of that
object.  Note that in this sense ``is a part of'' resembles ``is a
subset of,'' but does not resemble ``is an element of'' which is
neither reflexive nor transitive.
\begin{align*}
& \lforall[x][\Atom{\Obj P}{x,x}] \\
& \lforall[x][\lforall[y][((\Part{x}{y} \land \Part{y}{x})
      \lif \eq[x][y])]] \\
& \lforall[x][\lforall[y][\lforall[z][((\Part{x}{y} \land
        \Part{y}{z}) \lif \Part{x}{z})]]]\\
\intertext{Moreover, any two objects have a mereological sum (an object that has
  these two objects as parts, and is minimal in this respect).}  &
\lforall[x][\lforall[y][\lexists[z][\lforall[u][(\Part{z}{u} \liff
        (\Part{x}{u} \land \Part{y}{u}))]]]]
\end{align*}
These are only some of the basic principles of parthood considered by
metaphysicians.  Further principles, however, quickly become hard to
formulate or write down without first introducing some defined
relations.  For instance, most metaphysicians interested in mereology
also view the following as a valid principle: whenever an
object~$x$ has a proper part~$y$, it also has a part~$z$ that has no
parts in common with~$y$, and so that the fusion of $y$ and $z$ is
$x$.
\end{ex}





\olfileid{fol}{mat}{exr}

\olsection{Expressing Relations in \article{structure}
  \printtoken{S}{structure}}

\begin{explain}
One main use !!{formula}s can be put to is to express properties and
relations in !!a{structure}~$\Struct M$ in terms of the primitives of
the language~$\Lang L$ of~$\Struct M$.  By this we mean the following:
the !!{domain} of $\Struct M$ is a set of objects.  The !!{constant}s,
!!{function}s, and !!{predicate}s are interpreted in~$\Struct M$ by
some objects in~$\Domain M$, functions on~$\Domain M$, and relations
on~$\Domain M$.  For instance, if $\Obj A^2_0$ is in $\Lang L$, then
$\Struct M$ assigns to it a relation~$R = \Assign{{\Obj
  A^2_0}}{M}$. Then the formula $\Atom{\Obj A^2_0}{\Obj v_1, \Obj v_2}$
\emph{expresses} that very relation, in the following sense: if a
variable assignment~$s$ maps $\Obj v_1$ to $a \in \Domain{M}$ and
$\Obj v_2$ to $b \in \Domain M$, then
\[
Rab \text{\quad iff\quad} \Sat{M}{\Atom{\Obj A^2_0}{\Obj v_1, \Obj v_2}}[s].
\]
Note that we have to involve variable assignments here: we can't just
say ``$Rab$ iff $\Sat{M}{\Atom{\Obj A^2_0}{a, b}}$'' because $a$ and
  $b$ are not symbols of our language: they are !!{element}s
  of~$\Domain{M}$.

Since we don't just have atomic !!{formula}s, but can combine them
using the logical connectives and the quantifiers, more complex
!!{formula}s can define other relations which aren't directly built
into~$\Struct M$.  We're interested in how to do that, and
specifically, which relations we can define in !!a{structure}.
\end{explain}

\begin{defn}
Let $!A(\Obj v_1,\dots, \Obj v_n)$ be !!a{formula} of $\Lang L$ in
which only $\Obj v_1$,\dots, $\Obj v_n$ occur free, and let $\Struct
M$ be !!a{structure} for~$\Lang L$. $!A(\Obj v_1,\dots, \Obj v_n)$
\emph{expresses the relation}~$R \subseteq \Domain M^n$ iff
\[
Ra_1\dots a_n \text{\quad iff\quad} \Sat{M}{\Atom{!A}{\Obj
    v_1,\dots, \Obj v_n}}[s]
\]
for any variable assignment~$s$ with $s(\Obj v_i) = a_i$ ($i = 1,
\dots, n$).
\end{defn}

\begin{ex}
In the standard model of arithmetic~$\Struct N$, the !!{formula} $\Obj
v_1 < \Obj v_2 \lor \eq[\Obj v_1][\Obj v_2]$ expresses the $\le$
relation on~$\Nat$. The !!{formula} $\eq[\Obj v_2][\Obj v_1']$
expresses the successor relation, i.e., the relation $R \subseteq
\Nat^2$ where $Rnm$ holds if $m$ is the successor of~$n$. The formula
$\eq[\Obj v_1][\Obj v_2']$ expresses the predecessor relation.  The
!!{formula}s $\lexists[\Obj v_3][(\eq/[\Obj v_3][\Obj 0] \land
  \eq[\Obj v_2][(\Obj v_1 + \Obj v_3)])]$ and $\lexists[\Obj
  v_3][\eq[(\Obj v_1 + {\Obj v_3}')][v_2]]$ both express the $\Obj <$
relation.  This means that the predicate symbol~$<$ is actually
superfluous in the language of arithmetic; it can be defined.
\end{ex}

\begin{explain}
This idea is not just interesting in specific !!{structure}s, but
generally whenever we use a language to describe an intended model or
models, i.e., when we consider theories.  These theories often only
contain a few !!{predicate}s as basic symbols, but in the domain they
are used to describe often many other relations play an important
role.  If these other relations can be systematically expressed by the
relations that interpret the basic !!{predicate}s of the language, we
say we can \emph{define} them in the language.
\end{explain}

\begin{prob}
Find !!{formula}s in $\Lang L_A$ which define the following relations:
\begin{enumerate}
\item $n$ is between $i$ and $j$;
\item $n$ evenly divides $m$ (i.e., $m$ is a multiple of $n$);
\item $n$ is a prime number (i.e., no number other than $1$ and $n$ evenly
  divides~$n$).
\end{enumerate}
\end{prob}

\begin{prob}
Suppose the formula $!A(\Obj v_1, \Obj v_2)$ expresses the relation $R
\subseteq \Domain M^2$ in !!a{structure}~$\Struct M$.  Find formulas
that express the following relations:
\begin{enumerate}
\item the inverse $R^{-1}$ of $R$;
\item the relative product $R \mid R$;
\end{enumerate}
Can you find a way to express $R^+$, the transitive closure of~$R$?
\end{prob}

\begin{prob}
Let $\Lang{L}$ be the language containing a 2-place predicate symbol
$<$ only (no other !!{constant}s, !!{function}s or !!{predicate}s---
except of course~$\eq$). Let $\Struct{N}$ be the structure such that
$\Domain{N} = \Nat$, and $\Assign{<}{N} = \Setabs{\tuple{n,m}}{n <
  m}$. Prove the following:
\begin{enumerate}
\item $\{ 0 \}$ is definable in $\Struct{N}$;
\item $\{ 1 \}$ is definable in $\Struct{N}$;
\item $\{ 2 \}$ is definable in $\Struct{N}$;
\item for each $n \in \Nat$, the set $\{ n \}$ is definable in
  $\Struct{N}$;
\item every finite subset of $\Domain{N}$ is definable in
  $\Struct{N}$;
\item every co-finite subset of $\Domain{N}$ is definable in
  $\Struct{N}$ (where $X \subseteq \Nat$ is co-finite iff
  $\Nat \setminus X$ is finite).
\end{enumerate}
\end{prob}






\olfileid{fol}{mat}{set}

\olsection{The Theory of Sets}

Almost all of mathematics can be developed in the theory of sets.
Developing mathematics in this theory involves a number of things.
First, it requires a set of axioms for the relation~$\in$.  A number
of different axiom systems have been developed, sometimes with
conflicting properties of~$\in$.  The axiom system known as
$\Log{ZFC}$, Zermelo--Fraenkel set theory with the axiom of choice
stands out: it is by far the most widely used and studied, because it
turns out that its axioms suffice to prove almost all the things
mathematicians expect to be able to prove.  But before that can be
established, it first is necessary to make clear how we can even
\emph{express} all the things mathematicians would like to express.
For starters, the language contains no !!{constant}s or !!{function}s,
so it seems at first glance unclear that we can talk about particular
sets (such as $\emptyset$ or $\Nat$), can talk about operations on
sets (such as $X \cup Y$ and $\Pow{X}$), let alone other
constructions which involve things other than sets, such as relations
and functions.

To begin with, ``is an element of'' is not the only relation we are
interested in: ``is a subset of'' seems almost as important.  But we
can \emph{define} ``is a subset of'' in terms of ``is an element of.''
To do this, we have to find !!a{formula}~$!A(x, y)$ in
the language of set theory which is satisfied by a pair of
sets~$\tuple{X, Y}$ iff $X \subseteq Y$.  But $X$ is a subset of $Y$
just in case all !!{element}s of~$X$ are also !!{element}s of~$Y$.  So
we can define $\subseteq$ by the formula
\[
\lforall[z][(z \in x \lif z \in y)]
\]
Now, whenever we want to use the relation~$\subseteq$ in a formula, we
could instead use that formula (with $x$ and $y$ suitably replaced,
and the bound variable~$z$ renamed if necessary).  For instance,
extensionality of sets means that if any sets~$x$ and $y$ are
contained in each other, then $x$ and $y$ must be the same set. This
can be expressed by $\lforall[x][\lforall[y][((x \subseteq y \land y
    \subseteq x) \lif x = y)]]$, or, if we replace $\subseteq$ by the
above definition, by
\[
\lforall[x][\lforall[y][((\lforall[z][(z \in x \lif z \in y)] \land
    \lforall[z][(z \in y \lif z \in x)]) \lif x = y)]].
\]
This is in fact one of the axioms of $\Log{ZFC}$, the ``axiom of
extensionality.''

There is no !!{constant} for $\emptyset$, but we can express ``$x$ is
empty'' by $\lnot \lexists[y][y \in x]$.  Then ``$\emptyset$ exists''
becomes the !!{sentence}~$\lexists[x][\lnot \lexists[y][y \in
    x]]$. This is another axiom of~$\Log{ZFC}$.  (Note that the axiom
of extensionality implies that there is only one empty set.)  Whenever
we want to talk about $\emptyset$ in the language of set theory, we
would write this as ``there is a set that's empty and \dots'' As an
example, to express the fact that $\emptyset$ is a subset of every
set, we could write
\[
\lexists[x][(\lnot\lexists[y][y \in x] \land \lforall[z][x \subseteq
    z])]
\]
where, of course, $x \subseteq z$ would in turn have to be replaced by
its definition.

To talk about operations on sets, such as $X \cup Y$ and $\Pow{X}$, we
have to use a similar trick.  There are no function symbols in the
language of set theory, but we can express the functional relations $X
\cup Y = Z$ and $\Pow{X} = Y$ by
\begin{align*}
& \lforall[u][((u \in x \lor u \in y) \liff u \in z)]\\
& \lforall[u][(u \subseteq x \liff u \in y )]
\end{align*}
since the !!{element}s of $X \cup Y$ are exactly the sets that are
either !!{element}s of~$X$ or !!{element}s of~$Y$, and the
!!{element}s of $\Pow{X}$ are exactly the subsets of~$X$.  However,
this doesn't allow us to use $x \cup y$ or $\Pow{x}$ as if they were
terms: we can only use the entire !!{formula}s that define the
relations $X \cup Y = Z$ and $\Pow{X} = Y$. In fact, we do not know
that these relations are ever satisfied, i.e., we do not know that
unions and power sets always exist. For instance, the !!{sentence}
$\lforall[x][\lexists[y][\Pow{x} = y]]$ is another axiom
of~$\Log{ZFC}$ (the power set axiom).

Now what about talk of ordered pairs or functions?  Here we have to
explain how we can think of ordered pairs and functions as special
kinds of sets.  One way to define the ordered pair $\tuple{x, y}$ is
as the set $\{\{x\}, \{x, y\}\}$.  But like before, we cannot
introduce !!a{function} that names this set; we can only define the
relation $\tuple{x, y} = z$, i.e., $\{\{x\}, \{x, y\}\} = z$:
\[
\lforall[u][(u \in z \liff (\lforall[v][(v \in u \liff v = x)] \lor
  \lforall[v][(v \in u \liff (v = x \lor v = y))]))]
\]
This says that the !!{element}s~$u$ of~$z$ are exactly those sets which
either have $x$ as its only !!{element} or have $x$ and~$y$ as its
only !!{element}s (in other words, those sets that are either identical
to $\{x\}$ or identical to $\{x, y\}$).  Once we have this, we can say
further things, e.g., that $X \times Y = Z$:
\[
\lforall[z][(z \in Z \liff \lexists[x][\lexists[y][(x \in
        X \land y \in Y \land \tuple{x, y} = z)]])]
\]

A function $f \colon X \to Y$ can be thought of as the relation $f(x)
= y$, i.e., as the set of pairs~$\Setabs{\tuple{x,y}}{f(x) = y}$. We
can then say that a set~$f$ is a function from $X$ to $Y$ if (a) it is
a relation $\subseteq X \times Y$, (b) it is total, i.e., for all $x
\in X$ there is some $y \in Y$ such that $\tuple{x, y} \in f$ and (c)
it is functional, i.e., whenever $\tuple{x, y}, \tuple{x, y'} \in f$,
$y = y'$ (because values of functions must be unique). So ``$f$ is a
function from $X$ to $Y$'' can be written as:
\begin{align*}
 \lforall[u][(u \in f \lif {}] & \lexists[x][\lexists[y][(x \in X \land y \in
      Y \land \tuple{x, y} = u)]]) \land {}\\
 \lforall[x][(x \in X \lif {}] &
      (\lexists[y][(y \in Y \land \mathrm{maps}(f, x, y))] \land {}\\
& (\lforall[y][\lforall[y'][((\mathrm{maps}(f, x, y) \land
    \mathrm{maps}(f, x, y')) \lif y = y')]]))
\end{align*}
where $\mathrm{maps}(f, x, y)$ abbreviates $\lexists[v][(v \in f \land
  \tuple{x, y} = v)]$ (this !!{formula} expresses ``$f(x) = y$'').

It is now also not hard to express that $f\colon X \to Y$ is
!!{injective}, for instance:
\begin{multline*}
 f \colon X \to Y \land \lforall[x][\lforall[x'][((x \in X \land x' \in
     X \land {}]] \\
 \lexists[y][(\mathrm{maps}(f, x, y) \land \mathrm{maps}(f,
        x', y))]) \lif x = x')
\end{multline*}
A function~$f\colon X \to Y$ is !!{injective} iff, whenever $f$ maps $x, x'
\in X$ to a single~$y$, $x = x'$.  If we abbreviate this formula as
$\mathrm{inj}(f, X, Y)$, we're already in a position to state in the
language of set theory something as non-trivial as Cantor's theorem:
there is no !!{injective} function from $\Pow{X}$ to $X$:
\[
\lforall[X][\lforall[Y][(\Pow{X} = Y \lif
    \lnot\lexists[f][\mathrm{inj}(f, Y, X)])]]
\]

One might think that set theory requires another axiom that guarantees
the existence of a set for every defining property. If $!A(x)$ is a
formula of set theory with the variable~$x$ free, we can consider the
!!{sentence}
\[
\lexists[y][\lforall[x][(x \in y \liff !A(x))]].
\]
This !!{sentence} states that there is a set~$y$ whose !!{element}s
are all and only those $x$ that satisfy~$!A(x)$. This schema is called
the ``comprehension principle.''  It looks very useful; unfortunately
it is inconsistent.  Take $!A(x) \ident \lnot x \in x$, then the
comprehension principle states
\[
\lexists[y][\lforall[x][(x \in y \liff x \notin x)]],
\]
i.e., it states the existence of a set of all sets that are not
!!{element}s of themselves. No such set can exist---this is Russell's
Paradox.  $\Log{ZFC}$, in fact, contains a restricted---and
consistent---version of this principle, the separation principle:
\[
\lforall[z][\lexists[y][\lforall[x][(x \in y \liff (x \in z \land
      !A(x))]]].
\]

\begin{prob}
Show that the comprehension principle is inconsistent by giving
!!a{derivation} that shows
\[
\lexists[y][\lforall[x][(x \in y \liff x \notin x)]] \Proves \lfalse.
\]
It may help to first show $(A \lif \lnot A) \land (\lnot A \lif A)
\Proves \lfalse$.
\end{prob}




\olfileid{fol}{mat}{siz}

\olsection{Expressing the Size of \printtoken{P}{structure}}

\begin{explain}
There are some properties of structures we can express even without
using the non-logical symbols of a language.  For instance, there are
!!{sentence}s which are true in !!a{structure} iff the !!{domain} of
the !!{structure} has at least, at most, or exactly a certain
number~$n$ of !!{element}s.
\end{explain}

\begin{prop}
The !!{sentence}
\begin{multline*}
  !A_{\ge n} \ident \lexists[x_1][\lexists[x_2][\dots\lexists[x_n][{}]]]\\
  \begin{aligned}
  (\eq/[x_1][x_2] \land {}
  \eq/[x_1][x_3] \land \eq/[x_1][x_4] \land \dots \land \eq/[x_1][x_n] \land {}\\
\eq/[x_2][x_3] \land \eq/[x_2][x_4] \land \dots \land {} \eq/[x_2][x_n] \land {} \\
\vdots\\
\eq/[x_{n-1}][x_n])
\end{aligned}
\end{multline*}
is true in !!a{structure}~$\Struct M$ iff $\Domain M$ contains at
least $n$ !!{element}s. Consequently, $\Sat{M}{\lnot !A_{\ge n+1}}$ iff
$\Domain M$ contains at most~$n$ !!{element}s.

\end{prop}

\begin{prop}
The !!{sentence}
\begin{multline*}
!A_{= n} \ident \lexists[x_1][\lexists[x_2][\dots\lexists[x_n][{}]]] \\
\begin{aligned}
  (\eq/[x_1][x_2] \land {}
  \eq/[x_1][x_3] \land \eq/[x_1][x_4] \land \dots \land \eq/[x_1][x_n] \land {}\\
\eq/[x_2][x_3] \land \eq/[x_2][x_4] \land \dots \land {} \eq/[x_2][x_n] \land {} \\
\vdots\\
\eq/[x_{n-1}][x_n] \land {} \\
\lforall[y][(\eq[y][x_1] \lor \dots \lor \eq[y][x_n]]))
\end{aligned}
\end{multline*}
is true in !!a{structure}~$\Struct M$ iff $\Domain M$ contains
exactly $n$ !!{element}s.
\end{prop}

\begin{prop}
A !!{structure} is infinite iff it is a model of
\[
\{!A_{\ge 1}, !A_{\ge 2}, !A_{\ge 3}, \dots \}.
\]
\end{prop}

There is no single purely logical sentence which is true in~$\Struct
M$ iff $\Domain M$ is infinite.  However, one can give !!{sentence}s with
non-logical !!{predicate}s which only have infinite models (although
not every infinite !!{structure} is a model of them).  The property of
being a finite structure, and the property of being a
!!{nonenumerable} structure cannot even be expressed with an infinite
set of !!{sentence}s.  These facts follow from the compactness and
L\"owenheim--Skolem theorems.



\OLEndChapterHook





